% !TeX root = ../main.tex
\section{Benchmark of Available Models}
\label{sec:benchmark}

Published results use different test sets and metrics
(Section~\ref{sec:task}), so we ran the models with released weights under
one protocol. The aim is the use case of this review: electric guitar, in
mixes, written as tablature.

\subsection{Setup}
\emph{Test sets.} EGDB test clips (ids divisible by five) as DI and through
three amplifier models, plus its five real recordings~\cite{chen2022egdb};
EGSet12, twelve real electric performances through a miked amplifier with
string labels~\cite{pedroza2024tones}; the GOAT test split as DI and five
amplifier renders~\cite{loth2025goat}; the Guitar-TECHS musical excerpts as DI
and miked amplifier~\cite{pedroza2025guitartechs}; GuitarSet players 03--05,
accompaniment and solo~\cite{xi2018guitarset}; and band mixes built by placing
EGDB and Guitar-TECHS guitar over the drums, bass and vocals of MUSDB18
songs~\cite{rafii2017musdb18}. At most twelve items per group are used,
spread evenly; GOAT has seven per tone and EGDB real five. Guitar-TECHS and GOAT
labels were shifted per recording to remove a constant timing offset.

\emph{Models.} Note transcription: the GAPS model~\cite{riley2024gaps}; a
version of it fine-tuned by us on EGDB, Guitar-TECHS and GuitarSet players
00--02 in many amplifier tones (``tone-augmented CRNN''); Basic
Pitch~\cite{bittner2022basicpitch}; YourMT3+ (YPTF.MoE+Multi)~\cite{chang2024yourmt3};
and MuScriptor small, medium and large conditioned on the guitar classes (and
on bass, drums and voice for the mixes)~\cite{rouard2026muscriptor}. For the
multi-instrument models we score the notes of the three guitar classes.
Separation: \texttt{htdemucs\_6s}~\cite{rouard2023htdemucs} and a community
BS-RoFormer six-stem checkpoint run through MSST~\cite{lu2023bsroformer,solovyev2026msst}.
Fingering: a hand-position Viterbi with a string prior learned from
GuitarSet players 00--02. End-to-end: TabCNN~\cite{wiggins2019tabcnn} trained
on GuitarSet and its tone-augmented variant~\cite{pedroza2024tones}. Beats:
Beat This!~\cite{foscarin2024beatthis}. Everything ran on one NVIDIA GB10
(DGX Spark), except Basic Pitch and TabCNN, which ran on CPU.

\emph{Metrics.} F50 on raw notes (Section~\ref{sec:task}); Tab F1, an F50
match that is also on the reference string; for fingering alone, the share of
reference notes placed on the reference string.

\emph{Caveats.} The band mixes combine a guitar part with an unrelated song,
so guitar and backing share neither key nor tempo; models that use musical
context, such as the sequence-to-sequence transcribers, may be penalised.
Both TabCNN models were trained on GuitarSet, probably including players
03--05. MuScriptor's training set is private, so overlap with these test sets
cannot be ruled out. Our tone-augmented CRNN saw other clips and players of
EGDB, Guitar-TECHS and GuitarSet, but no GOAT, EGSet12 or MUSDB18 audio.

\subsection{Note transcription}
Table~\ref{tab:bench-notes} gives F50 per set. No model wins everywhere.
On amplified electric guitar the models trained on many tones lead:
MuScriptor large reaches 0.77 F50 on the GOAT amplifier renders, against 0.71
for the tone-augmented CRNN and 0.56 for the GAPS model, with the fewest
octave-low errors (2.1\% of notes, against 8.0\% for GAPS) and the best note
ends (offset F1 0.68). MuScriptor improves steadily with size on amplified
material (0.66, 0.76, 0.77). On EGSet12, real electric guitar through a clean
amplifier, the GAPS model, the tone-augmented CRNN and MuScriptor large tie
(0.88--0.89). YourMT3+ is the best model on acoustic GuitarSet (0.92--0.93,
offset F1 0.83--0.85) but labels much amplified guitar as other instruments
(0.39 on GOAT amplifier renders, 0.54 on EGSet12). Basic Pitch trails the
tone-augmented CRNN by about 0.1--0.2 on every set.

\begin{table*}[t]
\centering
\caption{Note F50 (mean per item) on held-out sets. MS: MuScriptor
conditioned on the guitar classes; multi-instrument models are scored on
their guitar notes. Speed: audio seconds per second on one GB10 GPU (Basic
Pitch on CPU).}
\label{tab:bench-notes}
\footnotesize
\begin{tabular}{@{}lccccccc@{}}
\toprule
Set & GAPS & Tone-aug.\ CRNN & Basic Pitch & YourMT3+ & MS small & MS medium & MS large \\
\midrule
EGDB, DI + 3 amp models & 0.74 & \textbf{0.84} & 0.67 & 0.71 & 0.73 & 0.74 & 0.76 \\
EGDB real recordings (5) & 0.89 & \textbf{0.90} & 0.69 & 0.83 & 0.71 & 0.79 & 0.59 \\
EGSet12, real, miked amp & \textbf{0.89} & 0.88 & 0.76 & 0.54 & 0.85 & 0.88 & \textbf{0.89} \\
GOAT DI & 0.82 & \textbf{0.84} & 0.74 & 0.81 & 0.74 & 0.83 & 0.82 \\
GOAT, 5 amp renders & 0.56 & 0.71 & 0.56 & 0.39 & 0.66 & 0.76 & \textbf{0.77} \\
Guitar-TECHS excerpts & \textbf{0.81} & 0.78 & 0.61 & 0.70 & 0.66 & 0.73 & 0.71 \\
GuitarSet comp & 0.87 & 0.89 & 0.75 & \textbf{0.92} & 0.77 & 0.77 & 0.77 \\
GuitarSet solo & 0.91 & 0.92 & 0.81 & \textbf{0.93} & 0.86 & 0.89 & 0.86 \\
\midrule
Speed (\(\times\) real time) & 34 & 89 & 165 & 14 & 6.4 & 4.0 & 1.3 \\
\bottomrule
\end{tabular}
\end{table*}

\subsection{Separation}
Table~\ref{tab:bench-sep} scores each separator by the F50 of a transcription
of its guitar stem. Separation is decisive for the guitar-specific models:
the tone-augmented CRNN goes from 0.50 on the raw mix to 0.74 with
\texttt{htdemucs\_6s} and 0.81 with BS-RoFormer, close to its 0.82 on the same
guitar parts without backing. BS-RoFormer runs at 4.8 times real time against 39 for
\texttt{htdemucs\_6s}. The multi-instrument models gain less: MuScriptor large
rises from 0.44 to 0.75, and YourMT3+ is no better on the stem than on the
mix. Contrary to expectation, both transcribe the raw mixes worse than the
guitar-specific CRNNs; the unrelated backing described above is a plausible
cause that this benchmark cannot exclude.

\begin{table}[t]
\centering
\caption{F50 of the guitar notes in band mixes (41 mixes), transcribed from
the raw mix or from a separated guitar stem.}
\label{tab:bench-sep}
\footnotesize
\setlength{\tabcolsep}{4pt}
\begin{tabular}{@{}lccc@{}}
\toprule
Note model & Raw mix & htdemucs\_6s & BS-RoFormer \\
\midrule
GAPS & \textbf{0.51} & 0.72 & \textbf{0.81} \\
Tone-augmented CRNN & 0.50 & \textbf{0.74} & \textbf{0.81} \\
Basic Pitch & 0.36 & 0.58 & 0.64 \\
YourMT3+ & \textbf{0.51} & 0.45 & 0.52 \\
MuScriptor large & 0.44 & 0.64 & 0.75 \\
\bottomrule
\end{tabular}
\end{table}

\subsection{Tablature: modular versus end-to-end}
Table~\ref{tab:bench-tab} compares Tab F1. The modular pipelines lead on
every electric-guitar and mix condition. With the same Viterbi fingering,
MuScriptor large gives the best tablature on EGSet12 (0.63, against 0.57 for
the CRNNs at equal F50), and the tone-augmented CRNN the best on EGDB and on
separated mixes (0.59). End-to-end TabCNN is limited by note detection on
electric guitar (F50 0.36--0.50 on EGDB) rather than by string choice. The
exception is GuitarSet solo lines, where TabCNN-GPFX scores 0.81 against
0.42--0.55 for the modular pipelines. The cause is the fingering stage: from
the true notes, the Viterbi puts 80\% of accompaniment notes on the
reference string but only 56\% of solo notes, because its costs favour low
positions while solos move along the neck. On EGSet12 it reaches 69\%.

\begin{table*}[t]
\centering
\caption{Tablature F1 (note match on the same string). Modular rows place the
transcribed notes with the Viterbi fingering. The last row is the share of
reference notes the Viterbi puts on the reference string when given the true
notes. EGDB amps: mean of three amplifier models; mix columns: EGDB clips in
band mixes.}
\label{tab:bench-tab}
\footnotesize
\setlength{\tabcolsep}{4pt}
\begin{tabular}{@{}lccccccc@{}}
\toprule
Method & EGDB DI & EGDB amps & EGSet12 & GuitarSet comp & GuitarSet solo & Raw mix & BS-RoFormer stem \\
\midrule
GAPS + Viterbi & 0.52 & 0.44 & 0.57 & 0.71 & 0.55 & 0.36 & 0.47 \\
Tone-aug.\ CRNN + Viterbi & \textbf{0.60} & \textbf{0.55} & 0.57 & 0.73 & 0.53 & 0.33 & \textbf{0.59} \\
Basic Pitch + Viterbi & 0.45 & 0.43 & 0.40 & 0.46 & 0.42 & 0.28 & 0.45 \\
YourMT3+ + Viterbi & 0.59 & 0.48 & 0.30 & \textbf{0.76} & 0.54 & \textbf{0.43} & 0.46 \\
MuScriptor large + Viterbi & 0.53 & 0.53 & \textbf{0.63} & 0.64 & 0.50 & 0.37 & 0.54 \\
\midrule
TabCNN (end-to-end) & 0.40 & 0.32 & 0.31 & 0.71 & 0.76 & 0.18 & 0.34 \\
TabCNN-GPFX (end-to-end) & 0.43 & 0.37 & 0.40 & 0.74 & \textbf{0.81} & 0.20 & 0.37 \\
\midrule
Viterbi on true notes (string acc.) & \multicolumn{2}{c}{0.74} & 0.69 & 0.80 & 0.56 & -- & -- \\
\bottomrule
\end{tabular}
\end{table*}

\subsection{Beats and Noise2Fret}
On GuitarSet players 03--05, the larger Beat This! checkpoint reaches a beat
F-measure of 0.97 and a downbeat F-measure of 0.94 on accompaniment
(0.92 and 0.92 for the small one), but every tracker, including a
dynamic-programming baseline~\cite{ellis2007beat}, falls to about 0.5 beat
F-measure on unaccompanied solo lines.

Noise2Fret's released checkpoints~\cite{simionato2026noise2fret} reproduce
the scale of its reported results under its own protocol: tablature F-measure
0.92 on GuitarSet player 00 with and without the auxiliary losses. That
protocol scores 100\,ms windows and keeps only windows that contain an onset,
so silent passages cannot produce false notes; the number is not comparable
with the note-level Tab F1 of Table~\ref{tab:bench-tab}, and applying the
model to whole recordings needs an onset stage it does not include.

\subsection{Summary}
For transcribing electric guitar from a song, the strongest configuration we
measured is modular: a band-split separator, a guitar note model trained on
many tones (the tone-augmented CRNN, or MuScriptor large when speed allows),
and a fingering stage. Separation recovers most of the gap between mixes and
isolated guitar; the remaining weakness is fingering on melodic lines, where
a frame-level end-to-end model trained on acoustic guitar does better than a
hand-tuned Viterbi.
